\section{\algname for Dilated and Grouped Convolution}
\label{sec:zconv-alg-ext}

\ralg{zconv-ext} presents the \algname algorithm extended to support dilated and grouped convolution.
As in \ralg{zconv}, padding is omitted for readability.
The following changes are introduced in the extended version to support dilated and grouped convolution: 
\begin{enumerate*}[label=(\roman*)] 
    \item An extra buffer of shape \texttt{[OH,FW$\times$IC/GR]} is allocated at Line~\ref{alg:zconv-ext:alloc}.
    \item A new loop over the number of groups is added in Line~\ref{alg:zconv-ext:gr}. 
    \item A copy operation from the image tensor to the extra buffer is added before the GEMM call. 
    This copy handles stride, groups, and dilation, resulting in a contiguous matrix $A$.
    \item The GEMM parameters are adjusted for the extra buffer and to handle groups in matrices $B$ and $C$.
\end{enumerate*}

\begin{algorithm}[t]
\caption{Extended \algname algorithm.}
\label{alg:zconv-ext}
\small
\begin{algorithmic}[1]
\Require
\begin{tabular}[t]{@{}ll@{}}
$\texttt{IC}_{gr} \gets \texttt{IC}/\texttt{GR}$ & $\texttt{OC}_{gr} \gets \texttt{OC}/\texttt{GR}$ \\
$image[\texttt{N},\texttt{IH},\texttt{IW},\texttt{IC}]$ & $filter[\texttt{FH},\texttt{FW},\texttt{IC}_{gr},\texttt{OC}]$ \\
$output[\texttt{N},\texttt{OH},\texttt{OW},\texttt{OC}]$ & $bias[\texttt{OC}]$
\end{tabular}

\For{$n \gets 0$ to $\texttt{N}-1$} \Comment{Parallel loop}
  \For{$ow \gets 0$ to $\texttt{OW}-1$} \Comment{Parallel loop}
    \State Initialize output tile (same as \ralg{zconv})
    \State Allocate $A$ with size $\texttt{OH} \cdot \texttt{FW} \cdot \texttt{IC}_{gr}$ \label{alg:zconv-ext:alloc}

    \For{$fh \gets 0$ to $\texttt{FH}-1$}
      \For{$gr \gets 0$ to $\texttt{GR}-1$} \label{alg:zconv-ext:gr}
      \State $startH \gets fh \cdot \texttt{DH},\;\; endW \gets startH + \texttt{OH} \cdot \texttt{SH}$
      \State $startW \gets ow \cdot \texttt{SW},\;\; endH \gets startW + \texttt{FH} \cdot \texttt{DW} $
        \State $A \gets image[n,startH : endH : \texttt{SH}, startW : endW : \texttt{DW},gr\cdot \texttt{IC}_{GR} : (gr+1)\cdot \texttt{IC}_{GR}]$
        \State $B \gets \text{addr}(filter[fh,0,0,gr \cdot \texttt{OC}_{gr}])$ \label{alg:zconv-ext:b}
        \State $C \gets \text{addr}(output[n,0,ow,gr \cdot \texttt{OC}_{gr}])$ \label{alg:zconv-ext:c}
        \State $layout \gets rowMajor$
        \State $transa \gets noTrans,\;\; transb \gets noTrans$
        \State $m \gets \texttt{OH},\;\; k \gets \texttt{FW} \cdot \texttt{IC}_{gr},\;\; n \gets \texttt{OC}_{gr}$
        \State $lda \gets k,\;\; ldb \gets \texttt{OC},\;\; ldc \gets \texttt{OW} \cdot \texttt{OC}$
        \State $\alpha \gets 1,\;\; \beta \gets 1$
        \State $C \gets \alpha A B + \beta C$\Comment{GEMM call}
      \EndFor
    \EndFor
  \EndFor
\EndFor
\end{algorithmic}
\end{algorithm}

The extra buffer $A$ (Line~\ref{alg:zconv-ext:alloc}) is necessary to support dilated and grouped convolution.
Recall that in a row-major GEMM, rows must be contiguous in memory, and in the original algorithm, matrix $A$ has a row size of $\texttt{FW}\times\texttt{IC}$.
Dilation and grouping introduce irregular memory access patterns that break this contiguity.
Dilation in the width dimension inserts offsets between the \texttt{FW} elements of a row, while grouping reduces the number of input channels accessed, such that only $\texttt{IC}/\texttt{GR}$ elements are accessed per row instead of the full \texttt{IC} dimension.

Unlike in the original algorithm, matrix $B$ points to a non-contiguous tile of the filter tensor in the grouped convolution case.
When the number of groups is greater than one, each GEMM call processes only a subset of the output channels dimension (\texttt{OC}).
Specifically, \texttt{OC/GR} per group, rather than the full output channel dimension.
As a result, the base pointers to matrices $B$ and $C$ include an offset to the starting position of the current group, as shown in Lines~\ref{alg:zconv-ext:b} and~\ref{alg:zconv-ext:c}, and the corresponding row offsets of these matrices are handled through the leading dimension parameters \texttt{ldb} and \texttt{ldc}.

\subsection{DeepLabV3+ Model Evaluation}
\label{sec:deeplabv3plus}

The first three models in \ref{fig:torch-multi} stand out for obtaining the highest speedups with \algname.
These bars correspond to the following models: \emph{deeplabv3\_mobilenet\_v3\_large}, \emph{deeplabv3\_resnet50}, and \emph{deeplabv3\_resnet101}.
DeepLabV3 is a segmentation model architecture that utilizes dilated convolution and supports distinct backbones, such as \emph{resnet} and \emph{mobilenet}~\cite{Chen2017}.
Dilated convolutional layers typically have longer execution times due to their poor memory access patterns, as the elements within an input patch are non-contiguous in memory.
In DeepLabV3 models, dilated layers also have large padding sizes.
\algname's approach is to gather the necessary input elements into a contiguous buffer before executing a GEMM call.
While Im2col has a similar approach, \algname minimizes the number of non-contiguous accesses because it only considers the width dilation in the copy to the new buffer.
The height dilation only affects the loop that iterates over the filter height (\texttt{FH}).
Combined with its efficient padding handling, as discussed in \rsec{libtorch}, \algname significantly improves performance for these layers.

\rtab{deeplabv3plus} extends the evaluation to the more recent \emph{DeepLabV3+} models~\cite{Chen2018}.
While Timm and TorchVision do not provide these models, their PyTorch implementations are available on GitHub~\cite{FangGithub}.
Unlike the previous end-to-end model evaluations, these models are evaluated following the usage of the repository that implements them.
Thus, each model is run on a validation set of the Pascal VOC 2012 dataset, which contains 1449 images~\cite{Everingham10}.
\rtab{deeplabv3plus} presents details on the evaluated \emph{DeepLabV3+} model backbones, the total inference time for multithreaded execution across all validation images, and the corresponding speedup obtained by PyTorch-\algname.
For a simpler integration, the weight layouts remained unchanged after converting the model to the channel-last format, requiring \algname to convert the weight tensor to $(FH,FW,IC,OC)$ before execution.
Despite this additional transformation overhead, \algname improves execution time across all models, further demonstrating its effectiveness for this architecture.

\begin{table}[t]
  \caption{Inference times on Pascal VOC 2012 (1449 images) with \emph{DeepLabV3+}.}
  \label{tab:deeplabv3plus}
  \small
  \centering
  \begin{tabular}{l|rrr}
      Backbone & PyTorch (s) & PyTorch-\algname (s) & Speedup \\ \hline
      \emph{Xception}    & 213 & 189 & \textbf{1.13} \\
      \emph{Mobilenet}   &  92 &  89 & \textbf{1.03} \\
      \emph{ResNet50}    & 269 & 216 & \textbf{1.25} \\
      \emph{ResNet101}   & 341 & 289 & \textbf{1.18} \\
      \emph{HRNetV2\_32} & 589 & 456 & \textbf{1.29} \\
      \emph{HRNetV2\_48} & 883 & 678 & \textbf{1.30} \\
    \end{tabular}
\end{table}

\section{Artifact}

Upon acceptance, the authors will packaged the following components into an archival format and make it publicly available with a DOI:
\begin{enumerate*}[label=(\arabic*)]
  \item The figures and tables from the paper, and the raw data used to generate them.
  \item Source code for the \algname implementations (standalone and PyTorch integration).
  \item Benchmarking and evaluation scripts to reproduce measurements.
  \item A Dockerfile that details how to build all components required for the evaluation environment.
\end{enumerate*}
Full reproduction instructions will be included in the artifact README. % TODO: add instruction for deeplabv3plus evaluation